% !TEX root = main.tex
%---- Page setup, section styling, and the commands main.tex is written in ----

%---- Colour -----------------------------------------------------------------
% Headings and rules stay black: this document gets printed as often as it gets
% read on screen, and a colour heading is the first thing a laser printer ruins.
% The single accent is reserved for links, where colour carries meaning.
\definecolor{cvlink}{HTML}{123A70}
\definecolor{cvmuted}{HTML}{4A4A4A}

%---- Page --------------------------------------------------------------------
\geometry{
  a4paper,
  left=1.8cm, right=1.8cm,
  top=1.15cm, bottom=1.45cm,
  footskip=0.8cm,
}

% A CV that runs to two pages must survive being printed and separated, so every
% page carries the name and an "of N" count.
\pagestyle{fancy}
\fancyhf{}
\renewcommand{\headrulewidth}{0pt}
\renewcommand{\footrulewidth}{0pt}
\fancyfoot[L]{\footnotesize\color{cvmuted}\name}
\fancyfoot[R]{\footnotesize\color{cvmuted}Page \thepage\ of \pageref{LastPage}}

% Coloured links, no boxes. The border box is what makes an unstyled LaTeX CV
% look unfinished when it is opened in a browser rather than printed.
\hypersetup{
  colorlinks=true,
  urlcolor=cvlink,
  linkcolor=cvlink,
  pdfborder={0 0 0},
}

% No single line of an entry left behind at a page break.
\widowpenalty=10000
\clubpenalty=10000

\urlstyle{same}
\raggedright
\setlength{\parindent}{0pt}
\setlength{\tabcolsep}{0in}

%---- Sections ----------------------------------------------------------------
% Small caps + a full-width rule. Loud enough to scan in two seconds, quiet
% enough not to compete with the entries underneath.
% NOTE the braces around the [after-code]. Without them LaTeX's optional-argument
% scanner stops at the FIRST ']', which is \titlerule's own, and the stray
% trailing ']' then lands in the preamble as text.
% The rule needs clearance below the baseline or it strikes through the
% descenders of "Research & Course Projects".
\titleformat{\section}{\large\scshape\bfseries\raggedright}{}{0em}{}[{\vspace{1.2pt}\titlerule[0.9pt]}]
\titlespacing*{\section}{0pt}{2.6mm}{1.3mm}

%---- Entry lists -------------------------------------------------------------
\newcommand{\cvListStart}{\begin{itemize}[leftmargin=*, labelsep=0mm, label={}, topsep=0pt, parsep=0pt, itemsep=0.8mm]}
\newcommand{\cvListEnd}{\end{itemize}\vspace{0.4mm}}

% Two-column entry headers. Both columns are p-columns rather than l/r with
% \extracolsep{\fill}: an `l' column cannot wrap, so a long project title runs
% straight into the tool name on the right instead of breaking.
\newlength{\cvRightCol}\setlength{\cvRightCol}{4.0cm}
\newlength{\cvColGap}\setlength{\cvColGap}{3mm}
% \Needspace stops a break between an entry header and the bullets that explain
% it; a heading stranded at the foot of a page reads as a formatting accident.
\newcommand{\cvHeaderRows}[4]{%
  \Needspace*{5\baselineskip}%
  \begin{tabular}[t]{@{}>{\raggedright\arraybackslash}p{\dimexpr\textwidth-\cvRightCol-\cvColGap\relax}@{\hspace{\cvColGap}}>{\raggedleft\arraybackslash}p{\cvRightCol}@{}}
    \textbf{#1} & \textit{\small #2} \\[0.2mm]
    \textit{\small #3} & {\small #4}
  \end{tabular}%
  \vspace{-1.4mm}%
}

% {institution}{location}{role}{dates}
%   Institution ...................... dates
%   role ......................... location
\newcommand{\cvEntry}[4]{\item\cvHeaderRows{#1}{#4}{#3}{#2}}

% {title}{context: course, or how it came about}{tools}{note, e.g. a grade}
\newcommand{\cvProject}[4]{\item\cvHeaderRows{#1}{#3}{#2}{#4}}

% A one-line entry: {lead}{detail}{date}
% p-columns, not `l': an honour with a sentence of context has to wrap, and an
% `l' column silently runs off the page instead.
\newlength{\cvDateCol}\setlength{\cvDateCol}{2.6cm}
\newcommand{\cvLine}[3]{%
\item
  \begin{tabular}[t]{@{}>{\raggedright\arraybackslash}p{\dimexpr\textwidth-\cvDateCol-2mm\relax}@{\hspace{2mm}}>{\raggedleft\arraybackslash}p{\cvDateCol}@{}}
    \textbf{#1}#2 & \textit{\small #3}
  \end{tabular}%
  \vspace{-1.4mm}%
}

%---- Bullets -----------------------------------------------------------------
\renewcommand{\labelitemi}{\raisebox{0.25ex}{\tiny$\bullet$}}

\newcommand{\cvBulletsStart}{\begin{justify}\begin{itemize}[leftmargin=3.4ex, rightmargin=0.5ex, labelsep=1.4mm, topsep=0.6mm, parsep=0pt, itemsep=0.3mm]\small}
\newcommand{\cvBulletsEnd}{\end{itemize}\end{justify}\vspace{-1.6mm}}

%---- Skills table ------------------------------------------------------------
% A fixed label column so the five categories line up; without it the eye has to
% re-find the start of each value.
%
% Deliberately `tabular` and not `tabularx`: tabularx absorbs its own body while
% looking for \end{tabularx}, so it cannot be opened in one macro and closed in
% another. The value column is therefore sized by hand.
\newlength{\cvSkillLabel}\setlength{\cvSkillLabel}{5.4cm}
\newlength{\cvSkillGap}\setlength{\cvSkillGap}{2.5mm}
\newcommand{\cvSkillsStart}{%
  \small
  \begin{tabular}{@{}>{\bfseries\raggedright\arraybackslash}p{\cvSkillLabel}@{\hspace{\cvSkillGap}}>{\raggedright\arraybackslash}p{\dimexpr\textwidth-\cvSkillLabel-\cvSkillGap\relax}@{}}%
}
\newcommand{\cvSkillsEnd}{\end{tabular}\vspace{0.8mm}}
\newcommand{\cvSkill}[2]{#1 & #2 \\[0.6mm]}

%---- Masthead ----------------------------------------------------------------
% Identity left, contact right, as in template.tex. Used once, in main.tex:
% tabularx cannot be split across a start/end macro pair (see the note above).
\newcolumntype{L}{>{\raggedright\arraybackslash}X}

%---- Compact two-column lists (languages, certifications) --------------------
\newcommand{\cvPlainListStart}{\begin{itemize}[leftmargin=3.4ex, labelsep=1.4mm, topsep=0.6mm, parsep=0pt, itemsep=0.4mm]\small}
\newcommand{\cvPlainListEnd}{\end{itemize}\vspace{0.4mm}}
